% !TeX program = xelatex
% !TeX encoding = UTF-8
% 本文件自包含；优先使用 macOS 宋体，其他系统回退到 Fandol；无外部图片或 bibliography 依赖。
% 编译：latexmk -xelatex -interaction=nonstopmode -halt-on-error A题第一问_模型与求解算法.tex
\documentclass[UTF8,a4paper,zihao=-4,fontset=none]{ctexart}
\IfFontExistsTF{Songti SC}{
  \setCJKmainfont{Songti SC}[BoldFont={Heiti SC}]
  \setCJKsansfont{Heiti SC}
  \setCJKmonofont{Songti SC}
}{
  \setCJKmainfont{FandolSong-Regular.otf}[BoldFont=FandolHei-Regular.otf,ItalicFont=FandolKai-Regular.otf]
  \setCJKsansfont{FandolHei-Regular.otf}
  \setCJKmonofont{FandolFang-Regular.otf}
}
\usepackage[a4paper,margin=2.4cm,headheight=15pt]{geometry}
\usepackage{amsmath,amssymb,mathtools,bm}
\usepackage{booktabs,longtable,tabularx,array}
\usepackage{enumitem,fancyhdr,xcolor}
\usepackage[hidelinks,unicode]{hyperref}
\usepackage{bookmark}
\usepackage{setspace}
\setstretch{1.18}
\setlength{\parskip}{0.25em}
\setlength{\emergencystretch}{2em}
\setlist[enumerate]{leftmargin=2em,itemsep=0.25em,topsep=0.3em}
\ctexset{section={format=\Large\bfseries},subsection={format=\large\bfseries},subsubsection={format=\normalsize\bfseries}}
\pagestyle{fancy}
\fancyhf{}
\fancyhead[L]{\small 药材预热阶段的径向模型与算法}
\fancyhead[R]{\small 推导与审查资料}
\fancyfoot[C]{\thepage}
\renewcommand{\headrulewidth}{0.3pt}
\numberwithin{equation}{section}
\newcommand{\dd}{\mathop{}\!\mathrm d}
\newcommand{\pd}{\partial}
\newcommand{\units}[1]{\,\mathrm{#1}}
\newcommand{\source}[1]{\textbf{依据与性质。}#1\par}
\newcommand{\reason}[1]{\textbf{推导理由。}#1\par}
\newcolumntype{Y}{>{\raggedright\arraybackslash}X}
\title{\bfseries 圆柱形药材预热平衡阶段的\\温度与含水率模型及求解算法\\[0.5em]\large A题第一问：物理推导、守恒离散与精度验收}
\author{}
\date{2026年9月11日}
\begin{document}
\maketitle
\begin{abstract}
针对圆柱形药材在预热平衡阶段的热风处理过程，在内部状态仅随径向位置和时间变化的前提下，从薄壳能量守恒与水质量守恒出发，建立常物性导热方程和含水率相关的非线性扩散方程。表面采用有限阻力的换热与有效传质条件，明确区分题目给定条件、物理定律和补充闭合假设。数值算法采用包含轴心及表面节点的有限体积法、后向欧拉时间推进和 Picard 非线性迭代，并独立检查非线性残差、整体收支和完整输出范围内的网格及时间步收敛。本文完整解释各项方程与离散关系的来源，供模型审查和论文写作使用；不提供尚未完成精度验收的数值解。现有外部审查支持离散结构，但表明原起始网格与步长尚不足以满足目标精度。
\end{abstract}
\noindent\textbf{关键词}\quad 干基含水率；径向导热；非线性扩散；有限体积法；非线性残差；数值验证
\clearpage
{\setstretch{1.0}\fontsize{9.2}{9.7}\selectfont\tableofcontents}
\clearpage

\section{题面依据、研究对象与证据边界}
\subsection{第一问要求与原文定位}
研究对象为2026年全国大学生数学建模竞赛A题《药材的烘干问题》的第一问。以下引用均来自本地题目PDF第1页的第一问、第3页的附录1与附录2；文件索引见文末。引用只用于说明题意，不把题面未写出的物理条件补称为题目事实。

\textbf{Q1：几何。}第一问写明“某中药材形状大致呈圆柱形，长为25 cm，半径为2 cm”。据此采用圆柱作为理想几何，令长度$L=0.25\units m$、半径$R=0.02\units m$。“大致”意味着圆柱是几何近似，并不保证材料均质或边界受热对称。

\textbf{Q2：初始状态。}第一问写明“烘干开始时，药材的温度为28°C，水分浓度（即干基含水率）为2.55 kg/kg”。据此取初始温度$T_0=28\,^\circ\mathrm C$、初始干基含水率$C_0=2.55$；将二者推广到内部所有位置，还使用了初始状态均匀的假设。

\textbf{Q3：外界输入。}第一问写明“烘房的温度（单位：°C）和水分浓度（单位：kg/kg）的变化情况见附件1”。附录1进一步说明附件给出各时间点的烘房温度和水分浓度，因此它们作为随时间变化的已知外界输入，而不是本问待求未知量。附件数据每隔$60\units s$给出一次；第一行数据为时间$0$、温度$28$、水分浓度$0.01963$。

\textbf{Q4：物性。}第一问要求建立“预热平衡阶段药材温度和水分浓度变化规律的数学模型（相关参数见附录2）”。附录2给出的常数和扩散经验式见表\ref{tab:params}。第一问不使用附录3、4中的变物性公式。

\textbf{Q5：结果。}题面要求给出100、300、600、900、1200、1500、1800 s及距中心0、0.5、1、1.5、2 cm处的结果，并要求“1800 s内每隔1 s、到药材中心距离每隔0.1 cm的完整结果”，保留四位小数。附件3的模板以整数秒从1开始，温度与水分浓度分别存入两个工作表。因此，完整验收范围必须包含第1秒的表面位置，不能只核查论文中的七个时刻。

\subsection{题意解释与本研究的范围}
将题面“到药材中心距离”解释为横截面内到圆柱轴线的径向距离$r$。理由是输出只给出一个从0到半径的空间坐标，没有轴向坐标；此外，本研究已经确定忽略不同高度的状态差异。输出形式支持此解释，但不能单独证明真实有限圆柱的端部影响严格为零。

本研究只保留$T=T(r,t)$和$C=C(r,t)$，定义域为$0\le r\le R$、$0\le t\le t_f$，其中$t_f=1800\units s$。温度使用摄氏度，含水率使用$\mathrm{kg}\text{水}/\mathrm{kg}\text{干物质}$；方程内温度只出现在导数或温差中，摄氏温差与开尔文温差数值相同。

本文的物理定律是对题目背景的建模选择，并非题面逐条给出的方程。特别是傅里叶导热定律、牛顿换热定律、有效菲克扩散定律分别给出通量与状态差的关系，能量守恒、水质量守恒则把这些通量连接到状态的变化。仅有参数名称，并不能替代对定律适用条件的说明。

\section{符号、物理概念与模型假设}
\subsection{温度、热量、质量比与通量}
温度表示冷热状态；热量表示因温差传递的能量。相同温度的两块材料不因温差产生净导热，因而导热的直接驱动力是温度梯度，而不是温度的绝对大小。比热容$c_p$表示单位质量材料升高单位温度所需的能量；在本问常物性近似下，体积为$\dd V$的材料升温$\dd T$需要的显热为$\rho c_p\,\dd V\,\dd T$。这里“显热”专指温度变化对应的能量，不包括汽化潜热。

干基含水率的定义为$C=m_w/m_d$，其中$m_w$是水质量、$m_d$是干物质质量。由$m_w=Cm_d$与总质量$m=m_d+m_w$直接得到水的湿基质量分数$m_w/m=C/(1+C)$。因此$C$可以大于1，不能把$C=2.55$解释成水占总质量的255\%。

定义$\rho_d$为单位总体积中的干物质质量，则局部水的体积质量浓度为$\rho_d C$。水质量通量$j_r$表示单位面积、单位时间通过的水质量，单位为$\mathrm{kg/(m^2\cdot s)}$。热流密度$q_r$表示单位面积、单位时间通过的热量，单位为$\mathrm{W/m^2}$。本研究统一规定沿半径增大的方向为正，因而向轴心流动的热流为负。

\begin{table}[htbp]
\centering
\caption{第一问参数及来源}\label{tab:params}
\small
\begin{tabularx}{\textwidth}{@{}l Y l@{}}
\toprule
符号 & 含义及数值 & 来源\\
\midrule
$R,L$ & 半径$0.02\units m$，长度$0.25\units m$ & Q1\\
$T_0,C_0$ & 初温$28\,^\circ\mathrm C$，初始干基含水率$2.55$ & Q2\\
$\rho$ & 固定密度参数$820\units{kg/m^3}$ & 附录2\\
$c_p$ & 比热容$2600\units{J/(kg\cdot K)}$ & 附录2\\
$k$ & 导热系数$0.36\units{W/(m\cdot K)}$ & 附录2\\
$h$ & 对流换热系数$25\units{W/(m^2\cdot K)}$ & 附录2\\
$h_m$ & 有效对流传质系数$8\times10^{-7}\units{m/s}$ & 附录2\\
$D(C)$ & $7\times10^{-9}\exp(-0.89/C)\units{m^2/s}$ & 附录2\\
$T_\infty,Y_\infty$ & 烘房温度、附件所列环境水分浓度 & Q3、附件1\\
$C_e$ & 与环境对应的等效平衡干基含水率 & 假设A5\\
$T_s,C_s$ & 表面值$T(R,t)$、$C(R,t)$ & 符号约定\\
$\rho_d$ & 单位总体积的干物质质量，水分方程中可约去 & 假设A2\\
\bottomrule
\end{tabularx}
\end{table}

\subsection{假设清单与不能由题面直接推出的内容}
\textbf{A0：径向状态与端面处理。}内部状态仅随$r,t$变化，这是本研究已确定的简化。为得到本文的纯径向模型，进一步忽略两个端面的热量与水分交换，只保留侧面交换。忽略端面传质是闭合这一简化的具体约定；若保留其总体贡献，即使不保留高度坐标，也需额外建立等效交换项。

\textbf{A1：均质连续介质及对称性。}药材近似为等效均质连续介质，侧面外界条件沿周向和轴向均匀，初温及初含水率在内部均匀。每个位置用一个温度代表局部材料与水分的共同温度，不分别建模各相温度。Q1提供几何依据，Q2提供初始数值，但均质、均匀与局部温度统一仍是补充假设。

\textbf{A2：固定几何与静止干物质骨架。}最初1800 s内忽略收缩、膨胀及干物质损失，令$R,L$固定，$\rho_d$在空间和时间上为常数。这使干物质不运动，水分变化只由水的相对迁移造成。Q4指定第一问的常物性参数，而收缩在后续问题中单独讨论，这支持分阶段简化；但题面没有证明第一问实际收缩恰好为零。

\textbf{A3：有效水分扩散。}把毛细迁移、组织内输运等微观过程合并为由含水率梯度驱动的有效扩散，使用附录2的$D(C)$，不另设压力驱动流动、重力输运或液水与水蒸气两个状态。局部含水率下降不被自动解释成当地水分汽化。

\textbf{A4：简化能量收支。}只保留侧面对流、内部导热和显热储存，忽略显式蒸发潜热、水分迁移携带的显热、辐射、接触传热及反应热。Q4给出的$k,h,\rho,c_p$支持这一可闭合基线，但“未给潜热参数”不能证明潜热效应很小。该假设可能影响温度预测，应作为物理限制保留。

\textbf{A5：有效传质边界。}约定$C_e(t)=Y_\infty(t)$，即把附件环境水分浓度直接视为有效平衡含水率，并把$h_m$解释为与此质量比差配套的有效系数。这不是严格的空气与药材相平衡定律；二者虽同以$\mathrm{kg/kg}$书写，分母通常分别为干空气和干药材。题目没有给出将空气湿度转换为药材平衡含水率的吸附关系，因此本文必须显式列出这一闭合假设。

\textbf{A6：外界过程已知。}烘房条件由附件给定，不因本根药材失水吸热而反向改变；相邻采样点之间作线性变化。这是使用附件数据的输入约定，而非新的材料经验式。

\subsection{密度参数的一致解释}
固定体积与干物质守恒意味着实际湿质量随含水率变化。如果同时把$\rho=820$解释为永远严格不变的实际湿密度，就会与持续失水冲突。因此，本文只将$\rho c_p$作为Q4指定的固定体积热容量近似，水质量守恒另用$\rho_d C$表达。不使用$\rho C$代替水体积浓度，也不在第一问中强行加入一个与固定参数冲突的变湿密度关系。

若另有依据把题给$\rho$解释为初始湿密度，才可通过初始总质量分解得到$\rho_d=\rho/(1+C_0)$。本基线的含水率方程与边界中$\rho_d$均可约去，故无需为状态求解作这一额外解释；估算绝对水质量或加入潜热时则需要明确它。

\section{从物理收支推导连续模型}
\subsection{圆柱薄壳与径向几何}
取任意$0<a<b\le R$，将半径$a$至$b$之间的材料作为控制体，即一段同心圆柱壳。圆柱面面积为$A(r)=2\pi rL$，体积元为$\dd V=2\pi rL\,\dd r$，壳体积为$\pi L(b^2-a^2)$。这些关系直接来自Q1采用的圆柱几何，而不是额外的传热定律。

控制体方法的核心是先计算穿过两个壳面的交换，再决定壳内状态如何变化。半径较大的界面面积较大，所以相同的单位面积通量会对应更大的总传递量。这正是圆柱方程不能直接使用平板二阶导数的原因。

\subsection{导热定律与能量守恒}
在连续均质材料中采用傅里叶定律$q_r=-kT_r$，下标$r$表示对半径求偏导，$t$下标表示对时间求偏导，例如$T_t=\partial T/\partial t$；两个重复$r$下标表示二阶径向导数。负号使热量沿温度降低的方向流动。当表面较热时$T_r>0$，于是$q_r<0$，热量向内传递。附录2给出$k$，而将该系数用于傅里叶定律属于材料导热的本构选择。

壳体能量的增加等于由内界面进入的热量减去由外界面离开的热量。A4排除了其他能量项，因此任意壳体满足
\begin{equation}
\frac{\dd}{\dd t}\int_a^b \rho c_pT\,2\pi rL\,\dd r
=2\pi aLq_r(a,t)-2\pi bLq_r(b,t).
\label{eq:heat-shell}
\end{equation}
左侧用温度相对于任意固定参考温度表达显热；参考常数不影响时间导数。右侧符号来自向外通量为正的约定，并非根据干燥时温度高低临时选择。

代入$q_r=-kT_r$，约去公共因子$2\pi L$，再利用微积分基本定理，右侧成为$[rkT_r]_a^b=\int_a^b\pd_r(rkT_r)\,\dd r$。由于区间$a,b$任意且材料性质按A1连续，积分恒等式对应局部能量方程
\begin{equation}
\rho c_pT_t=\frac1r\pd_r(rkT_r),\qquad 0<r<R.
\label{eq:heat}
\end{equation}
\reason{先使用能量守恒确定收支，再使用傅里叶定律确定通量，最后使用圆柱面积和微积分把整体收支局部化。左、右两侧单位均为$\mathrm{W/m^3}$。}

第一问$k$为常数，式\eqref{eq:heat}可展开成$T_t=\alpha(T_{rr}+T_r/r)$，其中$\alpha=k/(\rho c_p)$是热扩散率，单位为$\mathrm{m^2/s}$。这一变形仅由乘积求导和除以常数体积热容量得到，不适用于直接把变导热系数移出导数。

\subsection{有效扩散与水质量守恒}
采用干基含水率形式的有效菲克定律$j_r=-\rho_dD(C)C_r$。负号使水分从含水率高处向低处迁移；$\rho_d$将质量比梯度换算为水质量通量，使单位成为$\mathrm{kg/(m^2\cdot s)}$。附录2给出有效扩散系数，A2、A3提供固定骨架和等效梯度输运的前提。

壳内水质量为$\int_a^b\rho_dC\,2\pi rL\,\dd r$。根据水质量守恒，壳内水量变化等于两个界面的净流入，故
\begin{equation}
\frac{\dd}{\dd t}\int_a^b\rho_dC\,2\pi rL\,\dd r
=2\pi aLj_r(a,t)-2\pi bLj_r(b,t).
\label{eq:water-shell}
\end{equation}
本文将$C$视为有效总含水状态，没有独立内部汽化状态。水在材料内从液态转为气态也不会凭空消失；故不能在式\eqref{eq:water-shell}右端随意加入由$C_t$构造的“失水源项”。

代入有效菲克定律，利用常数$\rho_d$约去并使壳体区间任意，得到
\begin{equation}
C_t=\frac1r\pd_r\!\left[rD(C)C_r\right],
\qquad D(C)=D_0e^{-\gamma/C},\quad
D_0=7\times10^{-9}\units{m^2/s},\quad\gamma=0.89.
\label{eq:water}
\end{equation}
\reason{该推导与热方程使用相同的圆柱收支几何，但守恒对象是水质量。$D_0,\gamma$来自附录2；指数按题目以$\mathrm{kg/kg}$表示的含水率数值使用，$\gamma$与$C$取同一量纲约定，使指数无量纲。}

对$C>0$，链式法则给出$D'(C)=D(C)\gamma/C^2>0$，说明较低含水率对应较小扩散系数。乘积求导还给出$\frac1r\pd_r[rD(C)C_r]=D(C)(C_{rr}+C_r/r)+D'(C)C_r^2$。因此，把水分方程写成$C_t=D(C)(C_{rr}+C_r/r)$会遗漏一项；后续算法直接离散通量，避免这种错误。

\subsection{初始条件、轴心条件与其含义}
依据Q2的数值和A1的均匀初态假设，初始条件为$T(r,0)=T_0$、$C(r,0)=C_0$，对所有$0\le r\le R$成立。初始条件描述开始时的状态，不表示这些数值应在之后固定。

由于A0、A1的径向对称性和轴心光滑性，轴心附近不能偏向任意径向方向，因而$T_r(0,t)=0$、$C_r(0,t)=0$，并要求$T,C$有限。用平滑展开$u(r,t)=u(0,t)+\tfrac12u_{rr}(0,t)r^2+o(r^2)$可知$u_r/r\to u_{rr}(0,t)$，故圆柱算子在轴心取极限为$2u_{rr}(0,t)$。温度方程在轴心为$\rho c_pT_t=2kT_{rr}$，水分方程为$C_t=2D(C)C_{rr}$；后者的$D'(C)C_r^2$因对称性趋于零。

由此可见，$1/r$只是坐标表达中的形式奇点；中心状态仍能变化。数值上不应把$r$替换成任意小正数，也不应固定中心温度或含水率。

\subsection{表面对流换热与有效传质}
定义表面状态$T_s=T(R,t)$、$C_s=C(R,t)$。采用牛顿换热定律，空气向表面输入的热流密度为$h(T_\infty-T_s)$。在零厚度界面上不另设表面热容量，并按A4忽略潜热等项，则该供热等于传入内部的导热，得到
\begin{equation}
kT_r(R,t)=h[T_\infty(t)-T_s(t)].
\label{eq:heat-bc}
\end{equation}
\reason{内部向外的导热通量为$-kT_r$，向内为$kT_r$；将其与空气向内供热相等，便得到式\eqref{eq:heat-bc}。$h$来自附录2，温度输入来自Q3，具体的线性换热定律是本构选择。两侧单位均为$\mathrm{W/m^2}$。}

表面传质采用A5的经验闭合：向外水质量通量为$j_{\rm out}=\rho_dh_m(C_s-C_e)$。忽略独立的界面水分存储，内部到达表面的通量必须等于排向环境的通量；与有效菲克定律相等并约去$\rho_d$，得到
\begin{equation}
-D(C_s)C_r(R,t)=h_m[C_s(t)-C_e(t)],
\qquad C_e(t)=Y_\infty(t).
\label{eq:water-bc}
\end{equation}
\reason{式\eqref{eq:water-bc}中的左右两侧是已经约去$\rho_d$的通量，单位可写为$\mathrm{m/s}$乘以干基质量比，并非未经换算的水质量通量。它保持有限传质阻力；当$C_s>C_e$时$C_r(R,t)<0$，与表面较干的状态一致。}

式\eqref{eq:heat-bc}和式\eqref{eq:water-bc}称为Robin或第三类边界条件，即把表面值与法向导数联系起来。给出了有限$h,h_m$后，不能同时再强制$T_s=T_\infty$或$C_s=C_e$，否则对同一边界重复施加不相容约束。

\subsection{附件输入的连续化与启动相容性}
设附件相邻时间点为$\tau_\ell,\tau_{\ell+1}$，任一外界变量$z$表示$T_\infty$或$Y_\infty$。由A6的线性变化假设，对$\tau_\ell\le t\le\tau_{\ell+1}$定义
\begin{equation}
z(t)=(1-\theta)z_\ell+\theta z_{\ell+1},
\qquad \theta=\frac{t-\tau_\ell}{\tau_{\ell+1}-\tau_\ell}.
\label{eq:input}
\end{equation}
这是连接两个已知端点的直线公式，保证通过原始数据且不产生区间内超出两端点的新极值；它不宣称实际烘房在两次采样之间一定严格线性。第一问所需区间包含在附件范围内，不需要外推。

附件初始外温与药材初温相同，热边界与均匀初态相容。采用A5后，初始$C_e=0.01963$而$C_0=2.55$，均匀初态的$C_r=0$与瞬间要求的正向外通量不相容。因此水分边界对$t>0$施加，初态仍按Q2保留；起始瞬间可以缺少经典光滑性，随后形成表面梯度。早期边界层需要由细网格和小步长分辨，不能改动题给初值来消除它。

\subsection{完整定解问题与整体收支}
汇集式\eqref{eq:heat}--\eqref{eq:water-bc}，基线模型为
\begin{equation}
\left\{
\begin{aligned}
\rho c_pT_t&=\frac1r\pd_r(rkT_r),&
C_t&=\frac1r\pd_r[rD(C)C_r],&&0<r<R,\ t>0,\\
T(r,0)&=28,& C(r,0)&=2.55,&&0\le r\le R,\\
T_r(0,t)&=0,&C_r(0,t)&=0,&&t>0,\\
kT_r(R,t)&=h(T_\infty-T_s),&
-D(C_s)C_r(R,t)&=h_m(C_s-C_e),&&t>0.
\end{aligned}
\right.\,\text{。}
\label{eq:full}
\end{equation}
其中$D(C)$取式\eqref{eq:water}，外界输入取式\eqref{eq:input}，$C_e=Y_\infty$。该系统没有优化决策变量；给定输入后，待求对象为两个随空间和时间变化的场。由于$k,\rho,c_p$为常数、$D$只依赖$C$、热边界未保留潜热，两方程相互独立，不应称为双向耦合模型。

定义体积平均$\bar T=(2/R^2)\int_0^RT r\,\dd r$、$\bar C=(2/R^2)\int_0^RC r\,\dd r$。由$\dd V=2\pi Lr\,\dd r$得到这种带$r$权重的平均；简单径向算术平均不是体积平均。取$V=\pi R^2L$、$A_s=2\pi RL$，积分连续方程并代入边界，便得到
\begin{equation}
\rho c_pV\frac{\dd\bar T}{\dd t}=A_sh(T_\infty-T_s),
\qquad
\rho_dV\frac{\dd\bar C}{\dd t}=-A_s\rho_dh_m(C_s-C_e).
\label{eq:global}
\end{equation}
这与直接对整根圆柱使用能量及水质量守恒一致。$L$从局部状态方程消去，是内外交换面积与体积都含相同长度因子的结果，不表示长度对总储热量或总水质量没有影响。

\section{物理闭合的局限与可审查的扩展}
题目只给出了环境“水分浓度”的数值，没有给出空气含湿量与药材平衡含水率之间的吸附关系。若使用更完整的物理边界，应另给关系$C_e=\Phi(T_s,T_\infty,Y_\infty)$，或者从蒸气分压差建立传质通量。现有材料不足以识别任意未知函数$\Phi$，故本文保留A5，并把它与题目事实区分。

水蒸发所需汽化潜热不包含在普通比热容中。若将蒸发集中到表面，可保留内部方程而将热边界改成$kT_r=h(T_\infty-T_s)-L_vj_{\rm evap}$，其中$L_v$为汽化潜热，$j_{\rm evap}$为水质量通量。在A5的通量解释和向外蒸发条件下，可令$j_{\rm evap}=\rho_dh_m(C_s-C_e)$。乘积$L_vj_{\rm evap}$的单位为$\mathrm{W/m^2}$，表示供热中被相变消耗的部分。

此扩展不纳入本文基线算法。其采用需要补充$L_v$和$\rho_d$的取值依据，并审查表面平衡假设。不能将局部$-\rho_dC_t$全部当作当地汽化速率，因为含水率变化也可能只是内部水分重新分布。是否可忽略潜热，应比较相变耗热与供热、储热等项的量级；“预热阶段”这一名称不能代替检验。

模型参数、输入插值及边界解释的不确定性属于模型误差；之后的网格加密只减少数值误差。即使计算稳定到四位小数，也不能据此宣布真实药材状态已经达到四位小数的物理预测精度。

\section{有限体积离散：从连续收支到代数方程}
\subsection{方法选择与统一记号}
温度为常系数线性扩散，含水率为非线性扩散，二者空间上只需相邻位置交换信息。因此采用有限体积法直接离散控制体收支，采用后向欧拉处理时间，并用Picard迭代处理$D(C)$。有限体积法便于保持共享界面通量；后向欧拉对启动阶段不光滑的水分场较稳健，但其一阶时间误差仍必须通过实际加密检查。Picard迭代将一次非线性求解转成若干次正系数的三对角线性求解，不保证任意步长都收敛。

为统一两套推导，令$u$表示温度或含水率，连续方程写成$b u_t=r^{-1}\pd_r(ra u_r)$，侧面条件写成$a u_r(R,t)=s[u_e-u(R,t)]$。记号对应关系见表\ref{tab:unified}；水分边界只是将式\eqref{eq:water-bc}两侧乘以$-1$，物理方向没有改变。

\begin{table}[htbp]
\centering
\caption{离散推导使用的统一变量}\label{tab:unified}
\begin{tabular}{@{}llllll@{}}
\toprule
方程 & $u$ & $b$ & $a$ & $s$ & $u_e$\\
\midrule
温度 & $T$ & $\rho c_p$ & $k$ & $h$ & $T_\infty$\\
水分 & $C$ & $1$ & $D(C)$ & $h_m$ & $C_e$\\
\bottomrule
\end{tabular}
\end{table}

\subsection{节点、控制体和几何权重}
将$[0,R]$均匀划分为$N$段，节点$r_i=i\Delta r$、$\Delta r=R/N$，$i=0,\ldots,N$。两个端点均为未知节点，因而题目要求的轴心值和表面值可直接得到。令$N$为20的倍数，使$0.1\units{cm}=0.001\units m$间隔的21个输出位置全部落在节点上，输出索引为$i_j=jN/20$。

内部界面位于相邻节点中点$r_{i+1/2}=(r_i+r_{i+1})/2$。补充$r_{-1/2}=0$、$r_{N+1/2}=R$后，第$i$个控制体为$[r_{i-1/2},r_{i+1/2}]$，其几何权重为
\begin{equation}
w_i=\int_{r_{i-1/2}}^{r_{i+1/2}}r\,\dd r
=\frac{r_{i+1/2}^2-r_{i-1/2}^2}{2},
\qquad V_i=2\pi Lw_i.
\label{eq:weights}
\end{equation}
\reason{该权重是圆柱体积积分的精确几何因子；它并非平板中统一的$\Delta r$。轴心和表面控制体只有半个径向网格宽度。}

未知量$u_i(t)$近似节点处状态，用$w_i u_i$近似控制体内$\int u r\,\dd r$，即将控制体的存储量集中到节点。这是数值求积近似，不是连续方程恒等式；因此离散总量是连续总量的数值近似。内部通量取共享界面值，才能在相邻控制体间精确抵消。表面半控制体与初始边界层可能影响局部截断误差，不能仅凭中心差分形式就宣称所有位置均有相同的局部二阶误差；最终输出的空间阶由加密实测确认。

\subsection{界面通量与半离散方程}
在两个相邻节点间采用中心差商$u_r(r_{i+1/2})\approx(u_{i+1}-u_i)/\Delta r$。热问题的界面系数取$a_{i+1/2}=k$；水分问题取$a_{i+1/2}=\{D(C_i)+D(C_{i+1})\}/2$。算术平均保留正性，并在光滑系数条件下给出界面系数的一致近似；这里没有需要特殊界面条件的已知材料跃变。

定义$G_{i+1/2}=r_{i+1/2}a_{i+1/2}/\Delta r$。将统一连续方程乘以$r$并在第$i$个控制体积分，得到$b\,\dd(\int u r\,\dd r)/\dd t=[ra u_r]_{i-1/2}^{i+1/2}$。再使用上述节点存储量和界面差商近似，内部节点满足
\begin{equation}
bw_i\frac{\dd u_i}{\dd t}
=G_{i+1/2}(u_{i+1}-u_i)-G_{i-1/2}(u_i-u_{i-1}),
\qquad 1\le i\le N-1.
\label{eq:semidiscrete}
\end{equation}
方程右侧每个差值都表示对应相邻区域带来的净流入。正的$G$使较高状态向较低状态传递；这里直接离散$D(C)C_r$，自然包含连续式中变扩散系数的作用。

\subsection{后向欧拉时间推进}
设$t_{n+1}=t_n+\Delta t_n$，用$u_i^{n+1}$表示新时刻的未知状态。将时间导数近似为$(u_i^{n+1}-u_i^n)/\Delta t_n$，且右侧所有物理量取新时刻，式\eqref{eq:semidiscrete}变为
\begin{equation}
bw_i\frac{u_i^{n+1}-u_i^n}{\Delta t_n}
=G_{i+1/2}^{n+1}(u_{i+1}^{n+1}-u_i^{n+1})
-G_{i-1/2}^{n+1}(u_i^{n+1}-u_{i-1}^{n+1}).
\label{eq:be}
\end{equation}
\reason{后向差商来自时间导数的有限差分，新时刻右端使该方法为隐式方法，即必须同时解出新时刻各节点。对于光滑时间过程，其总体时间精度为一阶；初始不相容状态还需专门检查早期误差。输出每1秒一次是Q5的要求，不代表计算步长必须为1秒。}

\subsection{轴心控制体与表面控制体}
轴心内侧界面半径为零，界面总交换量为零，式\eqref{eq:be}去掉左侧不存在的通量后得到
\begin{equation}
bw_0\frac{u_0^{n+1}-u_0^n}{\Delta t_n}
=G_{1/2}^{n+1}(u_1^{n+1}-u_0^{n+1}).
\label{eq:center-discrete}
\end{equation}
由式\eqref{eq:weights}可算出$w_0=\Delta r^2/8$。若$a$为常数，则$G_{1/2}=a/2$，除以$w_0$后得到轴心空间项$4a(u_1-u_0)/\Delta r^2$。因为对称展开给出$u_1-u_0=\tfrac12u_{rr}(0)\Delta r^2+o(\Delta r^2)$，该项趋于$2a u_{rr}(0)$，与连续轴心极限一致。这是一项可直接人工核查的几何与系数检验。

表面控制体外界面使用已知的Robin通量$Ra u_r(R)=Rs(u_e-u_N)$，内界面仍与相邻控制体共享通量，所以
\begin{equation}
bw_N\frac{u_N^{n+1}-u_N^n}{\Delta t_n}
=Rs[u_e(t_{n+1})-u_N^{n+1}]
-G_{N-1/2}^{n+1}(u_N^{n+1}-u_{N-1}^{n+1}).
\label{eq:surface-discrete}
\end{equation}
\reason{表面节点代表的是一个有有限体积的半控制体，其存储项不能删除；连续界面本身仍没有独立存储。Robin条件已经通过外界面通量进入这一方程，不能再额外添加一个离散导数边界方程，否则重复约束表面未知量。}

\subsection{三对角系统的完整组装}
冻结水分扩散系数后，或直接对常系数温度求解时，令$\mu_i=bw_i/\Delta t_n$、$H=Rs$。内部节点的矩阵行由式\eqref{eq:be}移项得到
\begin{equation}
-G_{i-1/2}u_{i-1}^{n+1}
+(\mu_i+G_{i-1/2}+G_{i+1/2})u_i^{n+1}
-G_{i+1/2}u_{i+1}^{n+1}
=\mu_i u_i^n.
\label{eq:matrix-interior}
\end{equation}
轴心行与表面行分别由式\eqref{eq:center-discrete}、\eqref{eq:surface-discrete}移项得到
\begin{equation}
\begin{aligned}
(\mu_0+G_{1/2})u_0^{n+1}-G_{1/2}u_1^{n+1}&=\mu_0u_0^n,\\
-G_{N-1/2}u_{N-1}^{n+1}
+(\mu_N+G_{N-1/2}+H)u_N^{n+1}
&=\mu_Nu_N^n+Hu_e(t_{n+1}).
\end{aligned}
\label{eq:matrix-endpoints}
\end{equation}
上述三类行给出了全部$N+1$个未知量的$N+1$个方程。每行只含本节点与相邻节点，所以矩阵为三对角。$\mu_i>0$，扩散系数为正时$G>0$，从而对角元严格大于同一行非对角元绝对值之和。这个结构支持冻结系数线性系统的常规消元，并不等于非线性Picard迭代无条件收敛。

\section{线性求解、非线性迭代与时间步接受}
\subsection{温度线性系统与三对角追赶法}
温度的$b,a,s$均为常数，每个时间步只需一次三对角求解。若$\Delta t_n$不变，矩阵也不变，可以复用其分解，只更新旧时刻温度和外界输入对应的右端。

为说明“追赶法”不是额外假设，记任一三对角行为$\ell_i x_{i-1}+d_i x_i+v_i x_{i+1}=f_i$，其中$\ell_0=0$、$v_N=0$。第一行除以$d_0$后得到$x_0+\widehat v_0x_1=\widehat f_0$，其中$\widehat v_0=v_0/d_0$、$\widehat f_0=f_0/d_0$。若已由前一行得到$x_{i-1}=\widehat f_{i-1}-\widehat v_{i-1}x_i$，代入第$i$行即可消去$x_{i-1}$，得到
\begin{equation}
\delta_i=d_i-\ell_i\widehat v_{i-1},\qquad
\widehat v_i=\frac{v_i}{\delta_i},\qquad
\widehat f_i=\frac{f_i-\ell_i\widehat f_{i-1}}{\delta_i},
\quad i=1,\ldots,N.
\label{eq:thomas}
\end{equation}
最后一行给出$x_N=\widehat f_N$，然后从$i=N-1$到0依次回代$x_i=\widehat f_i-\widehat v_i x_{i+1}$。每行只作常数次运算，故时间和工作存储均为$O(N)$。程序仍应检查主元非零、结果有限以及原线性系统残差，避免组装或实现错误。

\subsection{水分Picard迭代的每一步}
水分的$G$依赖新时刻未知的$C^{n+1}$，不能只计算一次旧状态扩散系数后就当成完整隐式解。Picard迭代的思想是先猜测新状态，由猜测形成扩散系数，解线性系统，再以解更新猜测，直到满足原非线性方程。

每一时间步先令$C^{(0)}=C^n$。第$m$次迭代使用$D_i^{(m)}=D(C_i^{(m)})$与$G_{i+1/2}^{(m)}=r_{i+1/2}(D_i^{(m)}+D_{i+1}^{(m)})/(2\Delta r)$组装式\eqref{eq:matrix-interior}和式\eqref{eq:matrix-endpoints}，线性解记为$C^{(m+1)}$。这里旧时间层右端始终为$\mu_iC_i^n$，不能换成$\mu_iC_i^{(m)}$，因为内层迭代不是时间向前推进。

每次得到$C^{(m+1)}$后，重新用它计算$D$与$G$。定义原非线性残差$F_i(C)$为以下形式，所有$G(C)$均从被检查的同一个状态$C$计算。
\begin{equation}
\begin{aligned}
F_0(C)&=\mu_0(C_0-C_0^n)-G_{1/2}(C)(C_1-C_0),\\
F_i(C)&=\mu_i(C_i-C_i^n)
-G_{i+1/2}(C)(C_{i+1}-C_i)\\
&\quad+G_{i-1/2}(C)(C_i-C_{i-1}),\qquad 1\le i<N,\\
F_N(C)&=\mu_N(C_N-C_N^n)
-H[C_e(t_{n+1})-C_N]\\
&\quad+G_{N-1/2}(C)(C_N-C_{N-1}).
\end{aligned}
\label{eq:nonlinear-residual}
\end{equation}
\reason{这些式子分别把三个控制体方程移到左侧，因而$F(C)=0$才表示完整非线性隐式步成立。残差中的节点下标$C_0$表示轴心候选值；初始统一含水率在本节用$C_i^0=2.55$区分。若复用上一次冻结系数，只能检查辅助线性系统，无法检查真正要求解的非线性系统。}

\subsection{两项迭代停止条件}
以新状态重新组装的线性冻结矩阵对角元作为尺度，记为$A_{ii}(C)$；它不是完整非线性Jacobian对角元。暂取$\varepsilon_C=10^{-10}+10^{-9}\max_i|C_i^{(m+1)}|$，数值按$\mathrm{kg/kg}$的状态量计，要求
\begin{equation}
\max_i|C_i^{(m+1)}-C_i^{(m)}|\le\varepsilon_C,
\qquad
\max_i\frac{|F_i(C^{(m+1)})|}{A_{ii}(C^{(m+1)})}\le\varepsilon_C.
\label{eq:picard-stop}
\end{equation}
前一项检查固定点迭代是否仍在明显变化；后一项检查原非线性方程是否满足。用正的对角系数归一化，将残差转为含水率尺度，减少控制体大小对判据的直接影响。该容差属于算法约定，并非真实解误差上界，尤其在很小步长下还需通过容差收紧试验检查累计代数误差。

最后一个冻结线性系统的残差也需独立检查，温度和水分均可采用范数后向误差$\eta_{\rm lin}=\|Ax-f\|_\infty/(\|A\|_\infty\|x\|_\infty+\|f\|_\infty)$，零分母且零残差时定义为0。暂以$\eta_{\rm lin}\le10^{-12}$为双精度线性求解检查门槛；它只检验给定矩阵的求解，不替代式\eqref{eq:picard-stop}。

\subsection{范围保持的理由与异常处理}
对于冻结的正$G$，若一个新时刻节点取全局最大值，则各相邻差值在收支式中都不能向它提供正的内部净流入。如果它还高于旧时刻全部节点和新时刻外界值，左侧存储增量为正，而右侧不大于零，矛盾。对最小值作同样论证，得到新状态位于旧状态和本步外界值构成的区间内。这是离散最大值性质的直接收支证明。

从而可逐步检查温度是否位于初始温度及历史外温的范围内，含水率是否位于初始含水率及历史$C_e$的范围内。附件的正$C_e$与正初值还提供正的含水率下界；在此区间$D(C)>0$。程序仅允许与浮点残差量级相称的微小范围偏差，不应设定任意宽的“物理容差”。

出现明显越界、负含水率、非有限数或异常主元时，应拒绝接受该步并诊断，不把结果强行截断到合法范围。截断改变状态，却未增加相应界面通量，会破坏收支并掩盖错误。若正确组装的冻结线性解明显违反上述范围性质，不能只靠不断缩步掩盖实现问题。

\subsection{步长、失败限额和同步接受}
初值取$T_i^0=28$、$C_i^0=2.55$，工作精度采用双精度，整个求解过程中不取四位小数。基本步长采用1秒的二进制分数，便于精确到达整数输出时刻；每步终点还应对齐附件折点及$t_f$。外界值由式\eqref{eq:input}在新时刻计算，不从药材表面状态反推。

暂设每步Picard迭代上限50次。若未满足式\eqref{eq:picard-stop}，则拒绝本步、将步长减半，以同一个已接受旧状态重试。温度虽可独立求解，也不能在水分失败时先把全局时间推进；使用共同时间网格时应将两场一起接受或一起回退。温度矩阵在步长改变后重新分解。

基本步长和后续加密步长只是起始计划；缩步成功后可先保持较小步长，不额外引入自动增步。实际每步长度必须记录，网格比较时不能仅报告名义步长。若缩步已使时间网格不再按比例对应，则还需核对实际最大步长与启动段步长，不能机械套用名义加密比估计收敛阶。

暂以$10^{-6}\units s$作为最小步长保护值，低于它仍无法完成非线性收敛则停止并保存故障时刻、状态、残差和迭代次数，不生成合格结果。该值是防止无界重试的算法限额，不能解释为物理时间尺度或精度保证。

\section{守恒检查、外部审查结论与精度验收}
\subsection{离散整体收支与其检查范围}
将式\eqref{eq:be}、\eqref{eq:center-discrete}和\eqref{eq:surface-discrete}对全部节点求和，内部共享界面通量一正一负抵消。乘以$\Delta t_n$后得到
\begin{equation}
\begin{aligned}
\rho c_p\sum_{i=0}^Nw_i(T_i^{n+1}-T_i^n)
&=\Delta t_nRh[T_\infty(t_{n+1})-T_N^{n+1}],\\
\sum_{i=0}^Nw_i(C_i^{n+1}-C_i^n)
&=\Delta t_nRh_m[C_e(t_{n+1})-C_N^{n+1}].
\end{aligned}
\label{eq:discrete-global}
\end{equation}
两式分别约去了$2\pi L$和$2\pi L\rho_d$。定义$\mathcal B_u^n=b\sum_iw_i(u_i^{n+1}-u_i^n)-\Delta t_nRs(u_e^{n+1}-u_N^{n+1})$，即可记录单步有符号收支偏差、累计有符号偏差$\sum_n\mathcal B_u^n$以及累计绝对偏差$\sum_n|\mathcal B_u^n|$，后者防止正负误差抵消造成误判。

检查程序应从已接受的状态增量与边界输入重新计算这些量，而不只输出组装时保存的“已经相等”的两个中间量。比较时同时报告绝对偏差与使用总储存变化、总交换量归一化的偏差，避免在真实交换接近零时用不稳定的相对误差。

整体守恒能检查几何权重、共享界面、边界符号和实现收支的一致性，但不能检查Picard迭代是否充分收敛。对于任意一组冻结的$G$，内部通量在求和时都会抵消，因此即使该组$G$不对应最终候选含水率，式\eqref{eq:discrete-global}仍可能成立。原非线性残差的局部各分量也可能相互抵消，故必须独立使用式\eqref{eq:picard-stop}。

\subsection{审查报告的已知结果及其性质}
表\ref{tab:review}转录本研究收到的审查报告，不是本文重新运行所得。报告检查原起始配置$N=200,\Delta t=0.25\units s$，并指出含水率差异最大处为第1秒的药材表面。这些结果说明原起始配置未通过精度验收，不说明核心离散被证伪。

\begin{table}[htbp]
\centering
\caption{审查报告提供的加密差异（非本文新计算）}\label{tab:review}
\begin{tabular}{@{}lll@{}}
\toprule
比较 & 温度最大差$/^\circ\mathrm C$ & 含水率最大差$/(\mathrm{kg/kg})$\\
\midrule
时间步$0.5\to0.25\units s$ & $4.78\times10^{-4}$ & $1.03\times10^{-3}$\\
网格$200\to400$ & $1.47\times10^{-5}$ & $4.01\times10^{-3}$\\
\bottomrule
\end{tabular}
\end{table}

报告另提供一个单次Picard迭代反例：总质量收支残差约为$2\times10^{-22}$，但原非线性残差约为容差的345倍。该反例说明“守恒通过”不能推出“非线性收敛”；正确的时间步接受逻辑必须拒绝这一候选解。现有摘要未提供反例全部网格、步长和时刻信息，本文不宣称已独立复现其数值。

报告未发现轴心、表面、贝塞尔解析特例及制造解公式错误，实际测试支持约二阶空间收敛。因此继续保留当前算法结构，重点调整网格与时间步长，而不是无依据地更换物理模型或离散格式。

\subsection{必须覆盖的完整输出集合}
按照Q5及模板，定义验收集合
\begin{equation}
\mathcal O=\{(r_j,t_n):r_j=0.001j\units m,\ j=0,\ldots,20;\
t_n=n\units s,\ n=1,\ldots,1800\}.
\label{eq:output}
\end{equation}
每个场有$1800\times21$个验收位置，尤其包括最初几秒、轴心与表面。$t=0$初值另行检查。论文表格只是完整输出的子集，不能替代式\eqref{eq:output}的精度验收。

\subsection{分别控制空间、时间和代数误差}
对任一场$U$，令$U_{N,\Delta t}$表示对应配置的未舍入结果。固定空间网格比较时间步，固定时间步比较空间网格，从而定义
\begin{equation}
\begin{aligned}
E_t^U(N,\Delta t)&=\max_{\mathcal O}|U_{N,\Delta t}-U_{N,2\Delta t}|,\\
E_r^U(N,\Delta t)&=\max_{\mathcal O}|U_{N,\Delta t}-U_{N/2,\Delta t}|.
\end{aligned}
\label{eq:refinement}
\end{equation}
同一位置比较不需要空间插值，因为$N$为20的倍数；整数秒输出通过时间对齐取得。每个最大值均记录发生的时间、半径以及两级解的原始数值。另单独记录第1秒表面含水率，但仍保留完整集合最大值，因为继续加密可能改变最差位置。

原$N=200,\Delta t=0.25\units s$仅作粗网格诊断。下一候选配置可取$N=400,\Delta t=0.125\units s$，但不预先判定合格；随后视误差依次加倍$N$或减半$\Delta t$。先定位限制精度的方向再优先细化，找到候选合格配置后重新交叉检查另一方向。细空间网格会更清楚地分辨早期瞬态，所以不能直接继承粗空间网格上的时间精度结论。

为给原$10^{-5}$状态量精度目标留出余量，暂采用表\ref{tab:tolerance}的经验门槛。它是数值验收约定，不是题目直接给出的要求，也不等于严格真实误差界；题面只要求结果保留四位小数。

\begin{table}[htbp]
\centering
\caption{相邻细化层级差异的暂定门槛}\label{tab:tolerance}
\begin{tabular}{@{}lll@{}}
\toprule
检查量 & 温度 & 含水率\\
\midrule
$E_t^U$ & $4\times10^{-6}\,^\circ\mathrm C$ & $4\times10^{-6}\units{kg/kg}$\\
$E_r^U$ & $4\times10^{-6}\,^\circ\mathrm C$ & $4\times10^{-6}\units{kg/kg}$\\
\bottomrule
\end{tabular}
\end{table}

相邻差异小还可能来自误差抵消，故需要三个连续层级检查下降趋势。若另一方向误差已控制，按误差比例$e\propto h^p$可得观测阶$p_r=\log_2[E_r^U(N,\Delta t)/E_r^U(2N,\Delta t)]$，时间方向同理为$p_t=\log_2[E_t^U(N,\Delta t)/E_t^U(N,\Delta t/2)]$。比值中若有零或已接近浮点与代数误差底限，不再报告有意义的观测阶。

一阶时间与约二阶空间是待实际检查的趋势，不是跳过加密的理由。若差异不降、观测阶不稳定或最大误差位置迁移，应继续诊断；最终结果取实际细网格解，不将Richardson外推值混入守恒结果。候选配置还须将Picard容差收紧一个数量级，再比较完整输出；暂要求变化小于$10^{-7}\units{kg/kg}$，以使代数误差明显小于离散验收门槛，否则继续收紧并复核收敛。

最后检查两级结果的四位小数一致性。靠近舍入分界时，很小误差也可能改变末位，所以必须保存原始双精度值。若只有极少数末位不一致，需继续细化或明确记录末位不确定性，不能通过事后截断掩盖。即使全部末位一致，也只是数值稳定性证据，不是物理预测精度的证明。

\section{独立验证的构造及推导}
\subsection{均匀平衡与全局收支测试}
令初始状态恒等于外界恒定状态，则所有梯度及边界差值都为零，连续方程右端为零，故状态应保持不变。这一结论直接由模型代入获得。离散算法若产生非零演化，则存在虚假源项、边界或线性组装错误。

对于非均匀状态，式\eqref{eq:discrete-global}提供全局收支核验；但它与主离散共享守恒结构，独立性有限，不能单独作为正确性结论。为检查几何与非线性实现，还需以下与真实未知解无关的确切解测试。

\subsection{常系数圆柱解析模态}
把统一方程中的$a,b,s$均取常数，外界$u_e$取常数，设扰动$u-u_e=A\phi(r)e^{-\lambda t}$。代入$b u_t=a(u_{rr}+u_r/r)$，约去公共时间因子后得到$\phi''+\phi'/r+(\lambda b/a)\phi=0$。令$\lambda=(a/b)\mu^2/R^2$，则这是零阶贝塞尔方程，轴心有限的解可取$\phi(r)=J_0(\mu r/R)$，其中$J_0$为第一类零阶贝塞尔函数。

利用导数恒等式$J_0'(x)=-J_1(x)$，表面$a u_r=s(u_e-u_s)$变为$-aA(\mu/R)J_1(\mu)e^{-\lambda t}=-sAJ_0(\mu)e^{-\lambda t}$。消去公共因子得到模态参数条件，连同确切解可写为
\begin{equation}
\mu J_1(\mu)=\frac{sR}{a}J_0(\mu),\qquad
u(r,t)=u_e+A J_0(\mu r/R)
\exp\!\left[-\frac{a}{b}\frac{\mu^2t}{R^2}\right].
\label{eq:bessel}
\end{equation}
\reason{该测试从连续方程分离变量得到，未使用主算法的矩阵公式，因此可以独立检验圆柱几何、轴心极限、表面Robin符号及时间衰减。选取正根$\mu$，初值设为式\eqref{eq:bessel}在$t=0$的值；它不是题目原始均匀初值。}

温度测试可沿用原$k,\rho c_p,h$；水分测试则人为把$D(C)$固定为常数$D_*>0$，并选取足够小的振幅$A$保证含水率为正。确切解的求根误差应远小于被检验数值误差。该测试不覆盖变扩散系数的非线性，故还需制造解。

\subsection{含水率非线性制造解}
制造解是先指定一个光滑正函数，再反推让它严格满足方程的已知源项；它用于测试程序，不是用人为源项修改真实药材模型。取$x=r/R$，选$C^*(r,t)=A(t)+B(t)x^2>0$，例如$A(t)=1+0.1e^{-t/\tau}$、$B(t)=0.1e^{-t/\tau}$，其中$\tau>0$是测试时间尺度。

直接求导可得$C_t^*=A'+B'x^2$、$C_r^*=2Br/R^2$、$C_{rr}^*=2B/R^2$。代入已经展开的非线性扩散算子，得到$r^{-1}\pd_r[rD(C^*)C_r^*]=(4B/R^2)[D(C^*)+Bx^2D'(C^*)]$。因此，在测试方程$C_t=r^{-1}\pd_r[rD(C)C_r]+f$中，应使用源项
\begin{equation}
f(r,t)=A'(t)+B'(t)x^2
-\frac{4B(t)}{R^2}\bigl[D(C^*)+B(t)x^2D'(C^*)\bigr].
\label{eq:manufactured}
\end{equation}
令初值为$C^*(r,0)$，中心导数天然为零；从原表面条件$-D(C_s)C_r=h_m(C_s-C_e)$解出测试外界输入$C_e^*(t)=C^*(R,t)+D(C^*(R,t))C_r^*(R,t)/h_m$，就得到初值、源项、边界完全自洽的确切解。

程序在测试控制体右端加入$\int_{r_{i-1/2}}^{r_{i+1/2}}f(r,t)r\,\dd r$，用独立且足够准确的积分计算该量，使源项积分误差显著小于被检验离散误差。该测试特别检查$D'(C)C_r^2$所对应的非线性影响；若主算法误把$D$简单移出导数，应无法通过相应加密验证。测试用源项与外界值不得进入正式题目计算。

\subsection{单次Picard误接受回归测试}
在同题参数下故意限制为一次冻结系数求解，同时计算整体收支与式\eqref{eq:nonlinear-residual}。若出现收支通过而原非线性残差超限，正确程序必须判为未收敛。精确复现外部报告数值需补齐反例的配置；即使尚未获得该配置，接受逻辑本身也必须覆盖此情形。

上述测试分别针对不同错误来源：平衡测试检查虚假演化，解析模态检查线性圆柱结构，制造解检查非线性算子，反例测试检查停止逻辑。它们不能证明A0、A4、A5的真实物理误差很小；物理假设仍需实验、背景关系或替代物理模型的比较。

\section{可实施流程、资源限额与结果记录}
\subsection{完整计算流程}
\begin{enumerate}[label=\arabic*.]
\item 读取题给参数与附件1，检查时间严格递增、单位一致、数值有限、数据覆盖$[0,1800]$，按式\eqref{eq:input}构造外界输入，保留原始数据不变。
\item 先执行均匀平衡、解析模态、制造解及误接受逻辑测试，确认基本结构正确，再进行真实题目配置的试算。
\item 建立$N+1$个节点与式\eqref{eq:weights}的权重，取均匀初值；整数秒输出使用$i_j=jN/20$，不采用原模板中省略号代表真实网格。
\item 确定本步终点，计算新时刻外界条件；组装并求解温度三对角系统，保存候选温度。
\item 以旧时间层含水率启动Picard迭代，逐轮更新$D,G$并解线性系统，重新计算原非线性残差，直至两项停止条件同时成立或达到迭代上限。
\item 若迭代失败，回退两场候选状态并缩步；若出现结构性范围异常、非有限数或达到最小步长，则停止并记录诊断。未经全部接受检查，不更新时间。
\item 每个接受步检查范围、线性误差、原非线性残差、单步及累计收支；在整数秒保存两场的21个指定节点及必要诊断。
\item 到达1800 s后，开展空间与时间分离加密及交叉确认，检查完整输出最大差、位置和收敛趋势；候选合格配置再作非线性容差收紧试验。
\item 精度和验证均满足后，保留双精度原始结果，按模板生成温度、水分浓度两个工作表；题目要求的七个时刻与五个位置从同一份结果抽取，最后统一设置四位小数显示。
\end{enumerate}

\subsection{规模、计算预算与停止边界}
设时间步数为$M$、每步水分平均迭代次数为$K$，每次三对角求解为$O(N)$，总计算量为$O(MNK)$，工作内存为$O(N)$。只保存题目要求的输出时，两场共$2\times1800\times21$个数值，存储规模与内部迭代次数无关；无需保存所有Picard中间解。

候选$N=400,\Delta t=0.125\units s$在不缩步时对应14400个时间步，仍未证明足够精确。先作短时段运行估计单位步成本，再决定加密批次；暂以单次10分钟、首轮批次30分钟作为资源保护设置，实际可根据实测调整。资源限额不是精度豁免，达到限额时应保存检查点并报告验证未完成，不以当前最好结果冒充合格结果。

每次重要运行至少记录输入文件校验值、模型与算法版本、$N$、名义与实际时间步、迭代容差、最大迭代次数、拒绝步数、执行环境、开始结束时间、停止原因、状态范围、线性及非线性残差、收支偏差和加密比较。该算法确定性运行，无需随机种子。检查点保存最后接受时刻及两场完整内部节点，使失败后的恢复不依赖不完整候选步。

\subsection{验收结论的表述边界}
“程序正常结束”“离散整体守恒”“Picard已收敛”“数值误差足够小”和“模型准确表达真实物理”是不同层次的结论，必须分别给出依据。现有审查支持基本离散结构，原起始配置未通过精度目标；本文整理了完整可实施流程，尚未给出经过全部加密与验证的最终配置，也没有新增数值解。

\appendix
\section{依据与公式追溯索引}
\begin{longtable}{@{}>{\raggedright\arraybackslash}p{0.22\textwidth}>{\raggedright\arraybackslash}p{0.34\textwidth}>{\raggedright\arraybackslash}p{0.35\textwidth}@{}}
\caption{主要数学关系的来源与审查前提}\label{tab:trace}\\
\toprule
关系 & 直接来源 & 需要保留的前提\\
\midrule
\endfirsthead
\toprule
关系 & 直接来源 & 需要保留的前提\\
\midrule
\endhead
圆柱壳体积与面积 & Q1的圆柱几何及圆柱体积公式 & 规则圆柱近似，A0、A2\\
初值 & Q2的温度与干基含水率数值 & 初始空间均匀，A1\\
$q_r=-kT_r$ & 傅里叶导热定律，附录2的$k$ & 连续介质、局部温度，A1\\
式\eqref{eq:heat} & 式\eqref{eq:heat-shell}的能量守恒及圆柱几何 & 固定体积热容量、无其他热源，A4\\
$j_r=-\rho_dD C_r$ & 有效菲克扩散，附录2的$D(C)$ & 固定干骨架与有效梯度输运，A2、A3\\
式\eqref{eq:water} & 式\eqref{eq:water-shell}的水质量守恒 & 常数$\rho_d$，无干物质损失\\
式\eqref{eq:heat-bc} & 牛顿换热及界面能量平衡，附录2的$h$ & 无独立界面存储、不计潜热，A4\\
式\eqref{eq:water-bc} & 有效线性传质及界面水质量平衡 & $C_e=Y_\infty$为闭合假设，A5\\
式\eqref{eq:input} & Q3数据及两点间线性函数 & 采样间变化按直线近似，A6\\
式\eqref{eq:be} & 控制体积分、共享通量差商、后向欧拉 & 节点存储求积近似，不是物理定律\\
式\eqref{eq:matrix-endpoints} & 轴心零面积及表面Robin通量 & 两端各有半控制体，不重复约束\\
式\eqref{eq:picard-stop} & 非线性离散方程及迭代变化 & 必须重算最新状态的扩散系数\\
式\eqref{eq:discrete-global} & 离散界面通量相消 & 不足以证明非线性收敛\\
式\eqref{eq:refinement} & 同位置不同分辨率结果比较 & 完整输出、分离并交叉控制误差\\
式\eqref{eq:bessel} & 连续常系数方程分离变量 & 测试专用初值与常系数\\
式\eqref{eq:manufactured} & 指定光滑场后直接求导反推源项 & 源项和输入仅用于测试\\
\bottomrule
\end{longtable}

\section{资料来源与文件说明}
题面依据为《2026年高教社杯全国大学生数学建模竞赛题目：A题 药材的烘干问题》，第一问位于PDF第1页，附录1和附录2位于PDF第3页。源文件在本地资料目录\texttt{CUMCM2026Problems/A题/A题.pdf}。外界输入为同目录下\texttt{附件/附件1.xlsx}，输出格式依据为\texttt{附件/附件3/result1.xlsx}；上述路径均相对于数学建模工作目录。

算法阶段组织参考团队的\texttt{cumcm-subproblem-modeling}工作流。表\ref{tab:review}及单次Picard反例来自本研究收到的算法审查报告，报告摘要由使用者在对话中提供，未附独立可复现运行清单，因此本文将其作为外部审查证据而非本次新计算结果。物理方程在正文中从所述定律和假设自行推导，本文未据未查阅的教材编造引文、页码或实验依据。

本文件为自包含的XeLaTeX源文件，优先使用系统宋体，其他环境回退到TeX Live提供的Fandol字体，不依赖外部图片、数据读取或参考文献数据库。编译仅生成本文，不执行药材模型数值求解。
\end{document}
